% RESEARCH_DRAFT_SHA256: c7cc9ba57628e4e10024cbb449ae16795cb7b1cdcabbe862e932d33fc8e249f8
% Layout and source-reference wrappers expanded for the web renderer.
\documentclass[11pt]{article}
\usepackage[margin=25mm]{geometry}
\usepackage{fontspec,unicode-math}
\setmainfont{Latin Modern Roman}
\setmathfont{Latin Modern Math}
\usepackage{amsmath,amssymb,mathtools,xurl,hyperref}
\hypersetup{colorlinks=true,urlcolor=blue,linkcolor=blue}
\setlength{\parindent}{0pt}
\setlength{\parskip}{5pt}
\setlength{\emergencystretch}{3em}
\newcommand{\sref}[2]{\hyperlink{source:#1:#2}{[S#2]}}
\newcommand{\cataloguescope}[1]{\paragraph{Source question.}#1}
\newcommand{\floor}[1]{\left\lfloor#1\right\rfloor}
\begin{document}

\section*{A metacomplexity survey fragment: why incompressibility does not prove average-case hardness}
\noindent UnsolvedMath ID 30004924; OWR-8415359-014. Research attempt dated 11 October 2026.
\paragraph{Result in scope.}This manuscript contains a source-specific attempt and complete proofs of the stated auxiliary calculations. It makes no new solution or novelty claim for the entire target.
\section{Original problem and preserved definition}
The following definition and source list are reproduced from the repository. Any earlier inspection dates and reported literature claims in this preserved text describe that source record, not additional verification by this attempt.
% BEGIN PRESERVED DEFINITION
\subsection{Definitions and mathematical statement}

\paragraph{Shared notation.}Standard number systems are $\mathbb N,\mathbb Z,\mathbb Q,\mathbb R,\mathbb C$. Specialized notation and index ranges follow the cited source and the definitions below.


The minimum circuit-size problem asks whether a Boolean function supplied by its truth table has a circuit below a given size bound. Time-bounded Kolmogorov complexity measures the shortest description producing a string within the prescribed time bound. One-way functions are efficiently computable but difficult to invert on their designated input distribution.
\paragraph{Record scope and qualifications.}
The catalogue title is retained as metadata. Named certificates, contracts, bounds or frameworks in that title are not asserted as results of the original source extract.
The supplied registry flags this extraction as unrecoverable. It is displayed as a literal source fragment; the definitions provide context rather than a replacement formal statement.

The excerpt ends during a list of known characterizations and does not supply a complete new question or theorem.

Source-record limitation: Source flags the statement as unrecoverable: The linked report was inspected: this passage surveys proved characterizations of one-way functions and contains no explicit unresolved problem.
The source text below is rendered literally where necessary; damaged formulas and composite extracts are not silently repaired into a new mathematical claim.
The supplied catalogue marks this record as unrecoverable. Its extract and terminology are documented, but a complete original problem has not been reconstructed.
\paragraph{Statement recorded in the supplied source.}
“Meta-complexity” refers to the computational complexity of problems that are themselves about complexity, eg., the Minimum Circuit Size Problem MCSP and the problem Kpoly of computing the polynomial-time bounded Kolmogorov com- plexity of a string. In this plenary talk, I gave a brief overview of recent work on meta-complexity, with an emphasis on characterizations of one-way functions by average-case hard- ness of various meta-complexity problems [2, 4, 3, 1], including MCSP, Kpoly and the (uncomputable) problem of computing Kolmogorov complexity. The relevant characterizations are: (1) The breakthrough in [2] where the first characterization of one-way func- tions by average-case hardness of a natural problem is given. Connecting Meta-Complexity and Crypto Rahul Santhanam \href{https://doi.org/10.4171/owr/2021/54}{[S2]}
\paragraph{Source status.}
The supplied catalogue labels this source record open; that metadata was checked against the supplied snapshot on 2026-10-11. The label is preserved for provenance and is not an independent verification of present open status or mathematical correctness. The catalogue literature assessment is dated 2026-08-22.
\subsection{Short English statement}
The fragment reports connections between metacomplexity and one-way functions.
\subsection{Sources}
\begin{itemize}
\item[S1] Ulam AI and the UnsolvedMath Contributors, UnsolvedMath: A Curated Collection of Open Mathematics Problems, record 30004924 (OWR-8415359-014). Catalogue curation and metadata: CC BY 4.0. The source extract is retained; definitions, English summaries and qualifications are editorial additions.
\url{https://huggingface.co/datasets/ulamai/UnsolvedMath}.
\item[S2] Complexity Theory (2022), Oberwolfach Reports, DOI 10.4171/owr/2021/54, pp. 2955--3005
\url{https://doi.org/10.4171/owr/2021/54}.
\item[S3] Ilango--Ren--Santhanam, meta-complexity characterization
\url{https://eccc.weizmann.ac.il/report/2021/082/}.
\item Source-specific Oberwolfach Reports licences, including CC BY-SA 4.0 where applicable, are retained; see THIRD\_PARTY\_NOTICES.md.
\end{itemize}
% END PRESERVED DEFINITION
\subsection{Verbatim supplied catalogue question}
“Meta-complexity” refers to the computational complexity of problems that are themselves about complexity, eg., the Minimum Circuit Size Problem MCSP and the problem Kpoly of computing the polynomial-time bounded Kolmogorov com- plexity of a string. In this plenary talk, I gave a brief overview of recent work on meta-complexity, with an emphasis on characterizations of one-way functions by average-case hard- ness of various meta-complexity problems [2, 4, 3, 1], including MCSP, Kpoly and the (uncomputable) problem of computing Kolmogorov complexity. The relevant characterizations are: (1) The breakthrough in [2] where the first characterization of one-way func- tions by average-case hardness of a natural problem is given. Connecting Meta-Complexity and Crypto Rahul Santhanam

\section{The actual source passage}
The primary report, printed page2994, is a survey of established
characterizations, not an explicit new conjecture. The complete
passage lists different complexity notions and different
average-case distributions. The Ilango--Ren--Santhanam primary
abstract confirms that its statements concern approximation and
samplable or locally samplable distributions. We retain the
incomplete catalogue extract and do not invent a missing target.
Our mathematical attempt checks a tempting route from counting
incompressible strings to the claimed cryptographic hardness;
the route fails for an explicit reason.

\section{Counting short descriptions gives an easy classifier}
Fix a universal description machine and let $K(x)$ be plain
Kolmogorov complexity. There are $2^L-1$ binary programs of
length less than $L$, so at most that many distinct strings
have $K(x)<L$. For a uniform $n$-bit string and an integer $c$ with $0\le c\le n$,
\[
 \Pr\{K(x)<n-c\}\le\frac{2^{n-c}-1}{2^n}<2^{-c}.
\]
The same upper count applies to any time-bounded complexity,
because restricting running times can only decrease the number
of outputs of the short programs. Thus the constant algorithm
that always predicts $K(x)\ge n-c$ decides this particular
threshold predicate with error less than $2^{-c}$ under the
uniform distribution. For $c=\lceil3\log_2n\rceil\le n$,
its error is less than $n^{-3}$.

This does not make exact complexity easy, and does not refute a
carefully specified gap problem or sampler-based characterization.
It shows the precise failure of the attempted inference:
``most strings need long descriptions'' is a counting statement,
whereas a balanced or suitably distributed decision task may
require distinguishing rare structured strings. Counting alone
can instead make an unbalanced decision problem easy on average.

\section{The same obstruction for truth-table circuit complexity}
Consider a Boolean function on $m$ inputs, supplied by its
$N=2^m$-bit truth table. Count fan-in-two circuits with at most
$s\ge1$ gates over the sixteen binary Boolean operations.
Topologically order their gates. Each gate chooses its operation
and two predecessors among at most $m+s$ wires; the output
chooses one of those wires. Summing over gate counts at most
$s$ bounds the number of descriptions by
\[
 (s+1)(m+s)\bigl(16(m+s)^2\bigr)^s.
\]
Unused inputs and redundant descriptions only overcount functions,
which is harmless. Therefore a uniformly random truth table has
minimum circuit size at most $s$ with probability at most this
quantity divided by $2^N$. In regimes where its base-two
logarithm is at most $N/2$, always answering ``size exceeds $s$''
has error at most $2^{-N/2}$. Again a large typical circuit
size is not by itself average-case hardness of this threshold
decision task. This elementary counting argument does not prove
or disprove the source's gap-MCSP equivalence.

\section{Uncomputability and bounded time are distinct}
For completeness, exact unbounded $K$ is not computable.
If it were, for each $n$ choose the lexicographically first
$n$-bit string $x_n$ with $K(x_n)\ge n$; the short-program
count guarantees its existence, and computable $K$ would make
this search an algorithm. An algorithm containing a self-delimiting
encoding of $n$ followed by this fixed search program would
print $x_n$ using $2\lfloor\log_2n\rfloor+O(1)$ bits: encode
the binary length in unary and then the binary expansion.
For sufficiently large $n$ this is less than $n$, a contradiction.
The fixed program's simulation overhead changes only the constant
in plain description length.

By contrast, fix a computable time bound $t(n)$ large enough
for literal printing in the chosen machine. There is a constant
$c_0$ such that every $n$-bit string has a literal program of
length at most $n+c_0$ within this bound. Enumerate all programs
of length at most $n+c_0$, simulate each for $t(n)$ steps,
and select the shortest producing $x$. This finite procedure
computes $K^t(x)$ exactly. Its enumeration is exponential in
$n$; decidability is not a polynomial-time algorithm. If a time
bound does not allow the literal-printing premise, that premise
must be replaced by an explicit finite search convention, rather
than being silently assumed.

\section{Audit and remaining scope}
The audit checked the threshold's strict inequality, the
program-count union bound, the truth-table length versus input
arity, and the explicit printing-time assumption. The attempted
counting proof supplies no inversion reduction and no hardness
under the source's relevant sampler and approximation promises.
Known deep characterizations are attributed as context and not
re-proved or claimed as this attempt's result. No experiment or
theorem prover was used. Because the preserved fragment is a
survey rather than a complete open statement, its original
research target remains unreconstructed and completion is zero.

\section{Process and contributions}
Mathematics and model-assisted adversarial self-audit: OpenAI gpt-6-astra, ultra reasoning, as dispatched for this research batch; exact serving revision not exposed. Codex (GPT-6; exact serving revision not exposed) materially assisted with orchestration and format review. No human mathematical authorship, mathematical edits, or independent human proof checking is asserted.
The coordinating Codex (GPT-6; exact serving revision not exposed) also provided a separate model-assisted written-proof review of this manuscript. This is not independent human verification.
Contribution prompt A was applied to the full definition and supplied source row, then prompt B was applied adversarially to the derived results. The item-specific audit above records the scope checks and rejected extrapolations. Local assembly and structural checks do not establish mathematical correctness. The source catalogue status is preserved in TOP\_PROBLEM and is separate from this attempt status.
\paragraph{Integration review.} Codex (GPT-6; exact serving revision not exposed) reviewed the written arguments and expanded layout and source-reference wrappers. Repository metadata and the production web renderer were checked using Python 3.13 and JavaScript with MathJax 3.2.2. This model-assisted review is not independent human proof verification; no PDF compilation or formal proof checking is asserted.

\section{Completion Estimate}
\noindent\textbf{COMPLETION: 0\%}
The proved material is a rederivation of known or auxiliary facts rather than a substantiated new advance on the original target. The source extract surveys proved results and supplies no new complete question. Counting incompressibility does not establish the sampler- and gap-specific hardness needed for a one-way-function characterization.
\section{References}
The complete original Sources subsection above is retained, including catalogue attribution and licence notices. The additional list records primary sources actually consulted for this attempt on 11 October 2026; uncited source links reproduced above were not all rechecked.
\begin{itemize}
\item OWR54/2021 primary report; relevant complete contributions inspected from cached text, with printed pages specified in the manuscript. Historical known results are not treated as new contributions. \url{https://ems.press/content/serial-article-files/46932}.
\item Ilango, Ren and Santhanam, Hardness on any Samplable Distribution Suffices (ECCC2021/082); primary abstract inspected for the approximation and sampling quantifiers, without invoking its deep equivalence proof. \url{https://eccc.weizmann.ac.il/report/2021/082/}.
\end{itemize}
\section{Final status}
\textbf{Unresolved in this attempt.} The source extract surveys proved results and supplies no new complete question. Counting incompressibility does not establish the sampler- and gap-specific hardness needed for a one-way-function characterization.
\end{document}
